\section{模型的评价}

\subsection{模型的优点}

\begin{enumerate}
    \item \textbf{模型具有较好的普适性}。本文建立的功率平衡、离散调度、连续调节和离网储能模型均基于通用的电力系统运行原理和化工生产过程，不依赖于特定的设备参数或天气条件。模型中的装机容量、设备功率、电价等参数均可根据实际园区配置替换，能够适用于不同规模和不同地区的绿电直连电氢氨园区的优化运行分析。

    \item \textbf{优化算法高效且可保证全局最优性}。针对离散调度问题，采用边际成本排序法代替传统MILP求解器，将计算复杂度从指数级降至$O(24\log 24)$，且数学上可证明全局最优性。针对连续调节问题，采用贪心分配策略，24场景$\times$5产量的全枚举可在数秒内完成。

    \item \textbf{模型具有较好的鲁棒性}。通过24种风光场景（6风电$\times$4光伏）进行全面测试，覆盖了高风电低光伏、低风电高光伏、两者均低或均高等各种组合。模型在各类场景下均能稳定求解，结果符合物理规律——风光充裕场景成本低、匮乏场景成本高。

    \item \textbf{多维度指标综合评估}。不仅以吨氨成本作为经济性评价指标，还引入国家政策要求的三项绿电直连指标，从经济效益和政策合规两个维度综合评价各方案，更全面客观地反映不同方案的优劣。

    \item \textbf{图文结合的分析框架}。采用"模型建立—模型求解—图表展示—深入分析"的标准框架，共生成19张专业图表（功率平衡曲线、甘特图、热力图、Pareto曲线等），使复杂的数据关系一目了然。
\end{enumerate}

\subsection{模型的缺点}

\begin{enumerate}
    \item \textbf{时间分辨率有待提高}。本文以小时为时段单位，对风光出力的分钟级波动无法精确刻画。若采用15分钟或更短的时间分辨率，模型精度有望进一步提升，但计算复杂度也会相应增加。

    \item \textbf{未考虑设备启停损耗}。本文假设设备可自由启停且无额外成本，实际电解槽频繁启停会加速膜电极衰减。在目标函数中引入启停惩罚项后，调度方案可能更趋保守。

    \item \textbf{模型参数为静态设定}。分时电价、上网电价、设备效率等参数均为固定值，未考虑其随时间的变化。在更长期的规划分析中，引入参数的动态变化将使模型更具前瞻性。
\end{enumerate}

\subsection{改进方向}

未来研究可从以下方面深入：采用更高时间分辨率和更精细的设备模型，将氢储能和氨储运纳入统一优化框架，采用随机规划或鲁棒优化方法处理风光出力的不确定性\cite{turner2022}，以及结合碳排放交易市场将碳收益纳入经济性评估模型。
